% =====================================================================
%  2bat-ud1-7.tex  —  SESSIÓ 7: El producte de matrices: la idea
%  (sense preàmbul: s'inclou des de main.tex)
%  Criteri: C3 (producte de matrices i compatibilitat). CONCEPTE DELICAT.
% =====================================================================
\horatitol{Sessió 7 — El producte de matrices: la idea}%
{D'on surt cada número quan multipliquem dues taules: el producte fila per columna, pas a pas.}

\subsection{Idea clau}

Fins ara hem sumat, restat i escalat matrices, i totes aquestes operacions es
feien \emph{casella a casella}. Però hi ha preguntes que això no resol. Imagina
que saps \textbf{quantes sessions} de cada tipus necessita un usuari i
\textbf{quantes hores} dura cada tipus de sessió. Per saber les \emph{hores
totals} de dedicació, no n'hi ha prou de sumar ni d'escalar: cal \textbf{combinar}
les dues taules d'una manera nova. Aquesta combinació és el \textbf{producte de
matrices}.

La regla és \emph{fila per columna}: per obtenir un número del resultat, agafem
una \textbf{fila} de la primera matriu i una \textbf{columna} de la segona,
multipliquem els seus elements \emph{un a un} i \emph{sumem} els productes.

\begin{idea}
\textbf{El producte, pas a pas}
\begin{itemize}[leftmargin=1.4em, itemsep=2pt]
  \item Cada element del resultat surt de creuar una \textbf{fila} de la primera
        amb una \textbf{columna} de la segona.
  \item Es multiplica element a element (primer amb primer, segon amb segon\dots)
        i després se \emph{sumen} tots aquests productes.
  \item \textbf{Compatibilitat:} per poder fer-ho, la fila i la columna han de
        tenir la \emph{mateixa longitud}. És a dir, el nombre de \textbf{columnes}
        de la primera matriu ha de coincidir amb el nombre de \textbf{files} de la
        segona. (Ho treballarem a fons més endavant.)
\end{itemize}
\end{idea}

\subsection{Exemples}

\begin{exemple}[d'on surt un número del resultat]
Un usuari necessita, a la setmana, $2$ sessions individuals i $3$ sessions
grupals. Cada sessió individual dura $1$ hora i cada sessió grupal, $2$ hores.
Quantes hores de dedicació són, en total?

Ho posem com a fila (sessions) i columna (hores):
\[
\begin{pmatrix} 2 & 3 \end{pmatrix}\cdot\begin{pmatrix} 1 \\ 2 \end{pmatrix}
= 2\per 1 + 3\per 2 = 2 + 6 = 8 \text{ hores}.
\]
Multipliquem cada tipus de sessió per la seva durada i sumem: $8$ hores setmanals.
\end{exemple}

\begin{exemple}[un producte de dues matrices $2\times 2$]
Dos usuaris necessiten sessions individuals i grupals (files: usuaris; columnes:
tipus de sessió):
\[
S=\begin{pmatrix} 2 & 3 \\ 4 & 1 \end{pmatrix}.
\]
Cada tipus de sessió té una durada en hores i un cost en euros (files: tipus de
sessió; columnes: hores i cost):
\[
D=\begin{pmatrix} 1 & 20 \\ 2 & 15 \end{pmatrix}.
\]
El producte $S\cdot D$ dona, per a cada usuari, les hores totals i el cost total.
Per exemple, l'element fila 1, columna 1 (hores de l'usuari 1) surt de creuar la
\textbf{fila 1} de $S$ amb la \textbf{columna 1} de $D$:
\[
2\per 1 + 3\per 2 = 8 \text{ hores}.
\]
Fent-ho per a totes les posicions:
\[
S\cdot D=\begin{pmatrix} 2\per 1+3\per 2 & 2\per 20+3\per 15 \\ 4\per 1+1\per 2 & 4\per 20+1\per 15 \end{pmatrix}
        =\begin{pmatrix} 8 & 85 \\ 6 & 95 \end{pmatrix}.
\]
L'usuari 1 ocupa $8$ hores i costa $85\eur$; l'usuari 2, $6$ hores i $95\eur$.
\end{exemple}

\subsection{Exercicis resolts}

\begin{exresolt}[fila per columna, un número cada vegada]
Calcula el producte
\[
\begin{pmatrix} 3 & 2 \end{pmatrix}\cdot\begin{pmatrix} 4 \\ 5 \end{pmatrix}
\]
i explica què representaria si la fila fossin «$3$ àpats de migdia i $2$ de vespre»
i la columna «$4\eur$ el migdia i $5\eur$ el vespre».
\solucio
Creuem la fila amb la columna, multiplicant element a element i sumant:
\[
3\per 4 + 2\per 5 = 12 + 10 = 22.
\]
Representa el cost total: $3$ àpats a $4\eur$ més $2$ àpats a $5\eur$ fan $22\eur$.
\resposta{El producte és $22$; representa un cost total de $22\eur$.}
\end{exresolt}

\begin{exresolt}[un producte $2\times 2$ guiat]
Calcula $A\cdot B$ amb
\[
A=\begin{pmatrix} 1 & 2 \\ 3 & 0 \end{pmatrix},\qquad
B=\begin{pmatrix} 5 & 1 \\ 4 & 2 \end{pmatrix}.
\]
\solucio
Cada element surt de creuar una fila de $A$ amb una columna de $B$:
\begin{itemize}[leftmargin=1.4em, itemsep=2pt]
  \item Fila 1 $\times$ columna 1: $1\per 5 + 2\per 4 = 5 + 8 = 13$.
  \item Fila 1 $\times$ columna 2: $1\per 1 + 2\per 2 = 1 + 4 = 5$.
  \item Fila 2 $\times$ columna 1: $3\per 5 + 0\per 4 = 15 + 0 = 15$.
  \item Fila 2 $\times$ columna 2: $3\per 1 + 0\per 2 = 3 + 0 = 3$.
\end{itemize}
Col·loquem cada número a la seva posició:
\[
A\cdot B=\begin{pmatrix} 13 & 5 \\ 15 & 3 \end{pmatrix}.
\]
\resposta{$A\cdot B=\begin{pmatrix} 13 & 5 \\ 15 & 3 \end{pmatrix}$.}
\end{exresolt}

\subsection{Exercicis per fer}

\begin{exercicis}
\begin{exlist}
  \item Calcula els productes fila per columna:
        \begin{enumerate}[label=\alph*), leftmargin=1.6em, topsep=2pt]
          \item $\begin{pmatrix} 2 & 4 \end{pmatrix}\cdot\begin{pmatrix} 3 \\ 1 \end{pmatrix}$
          \item $\begin{pmatrix} 5 & 2 \end{pmatrix}\cdot\begin{pmatrix} 2 \\ 6 \end{pmatrix}$
        \end{enumerate}
  \item Un usuari fa $3$ sessions individuals (de $1$ hora) i $2$ grupals (de
        $2$ hores) cada setmana. Escriu-ho com un producte fila per columna i
        calcula les hores totals de dedicació.
  \item Calcula $A\cdot B$ amb
        $A=\begin{pmatrix} 2 & 1 \\ 0 & 3 \end{pmatrix}$ i
        $B=\begin{pmatrix} 4 & 2 \\ 1 & 5 \end{pmatrix}$.
        (Fes-ho posició a posició, com a l'exercici resolt 2.)
  \item Un centre atén dos usuaris amb sessions individuals i grupals:
        \[ S=\begin{pmatrix} 4 & 2 \\ 1 & 5 \end{pmatrix}, \]
        i cada tipus de sessió dura unes hores i té un cost:
        \[ D=\begin{pmatrix} 1 & 10 \\ 2 & 8 \end{pmatrix}. \]
        Calcula $S\cdot D$ i digues quantes hores i quin cost té cada usuari.
  \item Calcula $\begin{pmatrix} 1 & 2 & 3 \end{pmatrix}\cdot\begin{pmatrix} 2 \\ 1 \\ 4 \end{pmatrix}$.
        (Aquí la fila i la columna tenen $3$ elements: es multiplica i se sumen
        els tres productes.)
  \item \textbf{(Raonament)} Volem multiplicar una matriu de dimensió $2\times 3$
        per una altra. Mirant la regla de compatibilitat (les columnes de la
        primera han de coincidir amb les files de la segona), quantes files ha de
        tenir la segona matriu perquè el producte es pugui fer? Posa un exemple de
        dimensió vàlida per a la segona.
\end{exlist}
\end{exercicis}

% ---------------------------------------------------------------------
\fullsolucions{Sessió 7}
\begin{solucions}
\begin{sollist}
  \item (a) $2\per 3+4\per 1=6+4=10$ \quad (b) $5\per 2+2\per 6=10+12=22$
  \item $\begin{pmatrix} 3 & 2 \end{pmatrix}\cdot\begin{pmatrix} 1 \\ 2 \end{pmatrix}
        =3\per 1+2\per 2=3+4=7$ hores.
  \item $A\cdot B=\begin{pmatrix} 9 & 9 \\ 3 & 15 \end{pmatrix}$.
  \item $S\cdot D=\begin{pmatrix} 8 & 56 \\ 11 & 50 \end{pmatrix}$;
        \;l'usuari 1: $8$ hores i $56\eur$; l'usuari 2: $11$ hores i $50\eur$.
  \item $1\per 2+2\per 1+3\per 4=2+2+12=16$.
  \item La segona matriu ha de tenir $3$ files (per coincidir amb les $3$ columnes
        de la primera). Un exemple vàlid és una matriu $3\times 2$; el producte
        resultant seria $2\times 2$.
\end{sollist}
\end{solucions}
